\documentclass[ecta,nameyear,draft]{econsocart}
\usepackage{amsmath,amsfonts,amssymb,amsthm,booktabs,array,longtable}
\usepackage[T1]{fontenc}
\usepackage{mathpazo,microtype,needspace,xurl}
\usepackage{xr-hyper}
\RequirePackage[colorlinks,citecolor=blue,linkcolor=blue,urlcolor=blue]{hyperref}
\startlocaldefs
\newtheorem{theorem}{Theorem}
\newtheorem{proposition}[theorem]{Proposition}
\newtheorem{lemma}[theorem]{Lemma}
\newtheorem{corollary}[theorem]{Corollary}
\newcommand{\E}{\mathbb E}
\newcommand{\R}{\mathbb R}
\endlocaldefs
\setlength{\emergencystretch}{3em}
\makeatletter\def\copyright@text{Prepared for referee review}\makeatother

\setcounter{secnumdepth}{2}

\externaldocument{build/supp}[supp.pdf]
\begin{document}
\begin{frontmatter}
\title{Neural Bellman Operators}
\runtitle{Neural Bellman Operators}
\begin{aug}
\author[id=au1,addressref={add1}]{\fnms{Qian}~\snm{QI}\ead[label=e69main]{qiqian@pku.edu.cn}}
\address[id=add1]{Peking University}
\end{aug}
\begin{abstract}
This paper develops Neural Bellman Operators for constructive policy evaluation and feasible improvement under the original economic law. Action witnesses and recursive interval queries provide an explicit policy-loss account without assuming successful hidden-layer optimization. We implement a certificate-aware fixed-feature ReLU readout and compare it with a coefficient-matched Bernstein readout under common information and optimization budgets. Independent contexts test the associated work loss, and prospective reuse charges training, loading, verification and recovery. We also prove a Bellman-centered estimator whose expectation is the original policy-cost difference and whose range is controlled by local regret certificates. Its numerical and statistical errors are allocated separately. A continuous-law investment study and a two-control extension distinguish reliable return, achieved training accuracy, complete work and economic noninferiority. Historical adverse evidence remains part of the comparison.
\end{abstract}
\begin{keyword}
\kwd{Dynamic programming}\kwd{Neural approximation}\kwd{Policy certification}
\kwd{Information acquisition}\kwd{Computational economics}
\end{keyword}
\end{frontmatter}
\section{Introduction}\label{sec:introduction67}
A continuation value matters because it changes an economic decision. Its approximation error, the loss of the implemented policy, and the resources used to return that policy are different objects. A reliable numerical method must connect them without assuming that a fitted predictor is accurate or that a smaller loss certificate implies a better policy.

We develop Neural Bellman Operators along this chain. A constructive own-future continuation retains feasible action witnesses and admits policy-identical native and affine--ReLU realizations. Acquired-state feasibility and primitive regularity bounds connect its approximation budget to the original Bellman optimum. Recursive interval queries then evaluate only the needed states, without constructing a complete state lattice. A fitted selector proposes a query location; verified original-law endpoints and a safeguarded fallback supply the certificate. The fitted model cannot authorize an action by itself.

The first contribution of the present development is a training-to-verification contract. Strong convexity supplies a computable upper chord from already purchased endpoint bounds. The resulting pre-query screen licenses a rounded action before another Bellman query when it closes. A certificate-aware convex readout objective measures unresolved verification work rather than action-label error. We execute this objective with a constructive ReLU feature map, a coefficient-matched Bernstein control and a label-loss control, and report exact optimization gaps for the returned weights. Independent contexts test the work loss; recursively visited policy states are not treated as independent validation observations.

The second contribution connects numerical certification to economic comparison. The sum of optimal Bellman advantages along an actual policy trajectory has expectation equal to that policy's loss. Subtracting two such observations cancels the common initial optimal value. A verified local regret contract bounds the observation's range, which can be much smaller than the range of realized economic costs. We prove simultaneous cost-difference intervals that include the numerical evaluation error and give a prospective precision rule. The additional Bellman integrations are paid computations, not a policy-value oracle. The estimand remains the original expected policy-cost difference.

The third contribution is an explicit complete-work comparison. Cold services charge their construction, queries, path evaluation and output. Prospective reuse charges model loading, identity checks and recovery over a fixed sixty-four-block window. A two-control extension compares triangular refinements and trained two-output proposals under the same original-law certificate. Reliable return, achieved readout quality, economic noninferiority and a complete-work advantage are reported separately. Their coexistence is an empirical question, not a presumption of neural superiority.

These results build on classical dynamic programming and approximation rather than claiming priority for safe improvement or useful predictions. The constructive envelope uses Lipschitz extension \citep{mcshane1934} and separable distance transforms \citep{felzenszwalb2012}. Min-plus approximation and Bellman inequalities have established roles \citep{lakshminarayanan2014,wang2015}. Prediction-guided optimization and learned warm starts address related computational questions \citep{sakaue2023warm,sambharya2023}. Sequential value-based evaluation can reduce variance \citep{jiang2016}; our additional object is a paid original-optimum interval calculation with verified regret-scale support. The detailed theorem sections identify each inherited principle and the precise extra feasibility, coverage and arithmetic conditions used here.

The economic application retains the nonlinear investment primitives and continuous laws of the original NBO paper. The complete development companion reproduces the preceding article, including its adverse timings, unresolved economic contrasts and path-enclosure correction, and adds the full new derivations. Earlier controlled-diffusion, recursive-preference and game applications retain their own hypotheses in the unchanged historical companions. The new query and statistical conclusions are not silently extended to those distinct models.

\section{The economic operator and its policy account}\label{sec:core49}
Consider dates $t=0,\ldots,T$ and an invariant compact state domain $X\subset\R^d$. At date $t$, the action belongs to a nonempty compact set $A_t(x)$. The transition is $F_t(x,a,Z_{t+1})$, the current cost is $c_t(x,a)$, and the terminal cost is $g(x)$. Write $B_0=1$ and $B_t=\prod_{j<t}\beta_j$, with $\beta_j>0$. An implemented policy includes its observation rule, action selection and feasible numerical realization, not just its stored parameters.

For a continuation $v$, define
\begin{equation}\label{eq:operator49}
 Q_t(v;x,a)=c_t(x,a)+\beta_t\rho_t(v(F_t(x,a,Z))),\qquad
 \mathcal T_tv(x)=\inf_{a\in A_t(x)}Q_t(v;x,a).
\end{equation}
Throughout the deterministic comparison, the evaluations $\rho_t$ are monotone and cash invariant on a common domain containing all functions and constant translates used below. They are therefore nonexpansive in the supremum norm. Continuous objectives, effective compact feasible correspondences and measurable minimizers, or measurable arbitrarily near-minimizers with their error charged, are assumed. The stipulated conditional recursion gives Bellman comparison against all admissible adapted controls. In the experiments $\rho_t$ is conditional expectation, costs are bounded and innovations have their original continuous uniform law. The direct-cost inference in Section~\ref{sec:decisions49} is specific to expectation; it is not automatically an inference theorem for nonlinear recursive utility or games.

A Neural Bellman Operator is a construction and evaluation procedure for \eqref{eq:operator49}. It forms targets using its \emph{own} fitted $f_{t+1}$, returns a neural continuation and a feasible policy, and records an economic error allowance. It does not receive the optimal continuation as a training label. Arbitrary trained networks can be evaluated by the first result below. The subsequent constructive backend specifies how a network and its actor are obtained without assuming successful nonconvex optimization.

\subsection{Action-dependent errors and complete-policy loss}
Let $H(a)$ be a true current cost and $\widehat H(a)$ its continuation approximation. If $\widehat a$ is $\delta$-optimal for $\widehat H$ on the same feasible action set, then
\begin{equation}\label{eq:center49}
 H(\widehat a)-\inf_aH(a)\le\delta+
 \operatorname{osc}_{a}\{\widehat H(a)-H(a)\}.
\end{equation}
Indeed, compare $\widehat a$ with an arbitrarily near-optimal true action and cancel the approximate objectives. The infimum limit gives the result. Additive continuation levels cancel; action-dependent discrepancies do not. Under a change in future policies, the continuation must be transported to the new economic problem before this inequality licenses reuse. The full development edition retains the original quantitative transport, recursive-domain and economic-application results. No transport bound is presumed valid after an unaccounted change of constraints or preferences.

\begin{proposition}[One-sided policy sandwich]\label{prop:sandwich49}
Suppose, uniformly in state,
\[
 f_t-\mathcal T_tf_{t+1}\le u_t,\qquad
 \mathcal T_t^\pi f_{t+1}-f_t\le v_t,
 \qquad \ell_T\le f_T-g\le u_T.
\]
Then the actual feasible policy satisfies
\begin{equation}\label{eq:sandwich49}
 0\le J_0^\pi-V_0^*\le
 \sum_{t<T}B_t(u_t+v_t)+B_T(u_T-\ell_T).
\end{equation}
The bound is invariant to date-specific constant shifts of the fitted continuations, provided the associated residual accounts are shifted consistently.
\end{proposition}
\noindent The proof is in Supplement, Section~\ref{supp:proof-core49-1}.


A sharper valid residual account for the same continuation and the same deployed actor sharpens knowledge of the policy. It does not retrain that policy or change its realized cost. Conversely, subtracting two upper bounds in \eqref{eq:sandwich49} does not identify the difference in actual costs. These distinctions govern both the historical comparisons and the new experiments.

\section{Constructing and compiling a feasible neural policy}\label{sec:construction49}
Assume the primitive objective has, for feasible state--action pairs, state modulus $A_t^Q=C_{x,t}+\beta_tF_{x,t}L_{t+1}$ and action modulus $D_t=C_{a,t}+\beta_tF_{a,t}L_{t+1}$. Let
\[
 d_H(A_t(x),A_t(z))\le K_t^A\|x-z\|_1,
 \qquad L_t=A_t^Q+D_tK_t^A,\qquad L_T=L_g.
\]
The Hausdorff condition and the primitive moduli imply that $\mathcal T_tf_{t+1}$ is $L_t$-Lipschitz whenever $f_{t+1}$ is $L_{t+1}$-Lipschitz. All moduli and numerical integration procedures are supplied by the economic primitives, not fitted from successful outcomes.

At state nodes $x_{t,i}$ with covering radius $h_t$, form effective feasible $k_t$-nets. Evaluate their own-future objectives with certified error at most $e_t$. Minimize the computed labels and retain the corresponding \emph{original} action $a_{t,i}$. A Borel repair $\Gamma_{t,i}(x)$ is feasible at $x$ and satisfies
\begin{equation}\label{eq:repair49}
 \|\Gamma_{t,i}(x)-a_{t,i}\|_1\le K_t^A\|x-x_{t,i}\|_1+\zeta_t.
\end{equation}
The allowance $\zeta_t$ concerns displacement; feasibility itself may not fail. Define
\begin{equation}\label{eq:cone49}
 f_t(x)=\min_i\{y_{t,i}+L_t\|x-x_{t,i}\|_1\},\quad
 i_t(x)=\min\arg\min_i\{y_{t,i}+L_t\|x-x_{t,i}\|_1\}.
\end{equation}
Return $\pi_t(x)=\Gamma_{t,i_t(x)}(x)$. Terminal labels approximate $g$ with error $e_T$ and use the same cone construction with modulus $L_g$.

\begin{theorem}[Own-future construction with feasible witnesses]\label{thm:construct49}
The functions in \eqref{eq:cone49} have exact affine--ReLU realizations and Lipschitz moduli $L_t$ independent of label error. The returned policy is feasible and measurable. It satisfies
\begin{align}\label{eq:enclosures49}
 -e_t\le f_t-\mathcal T_tf_{t+1}&\le D_tk_t+e_t+2L_th_t,\\
 \mathcal T_t^\pi f_{t+1}-f_t&\le e_t+D_t\zeta_t,\notag
\end{align}
and hence
\begin{equation}\label{eq:constructionloss49}
 J_0^\pi-V_0^*\le
 \sum_{t<T}B_t(D_tk_t+2e_t+2L_th_t+D_t\zeta_t)
 +B_T(2e_T+2L_gh_T).
\end{equation}
With effective covers, repairs and certified queries at the prescribed accuracy, allocating each nonnegative term a share of a positive economic tolerance gives a terminating finite construction. A fixed numerical precision or a capped finite experimental ladder need not contain that allocation.
\end{theorem}
\noindent The proof is in Supplement, Section~\ref{supp:proof-core49-2}.


The Lipschitz-extension principle is classical \citep{mcshane1934}. A native min-plus envelope and its exact ReLU circuit, using the same labels, witnesses and ties, are \emph{one method with two encodings}. They return the same function and policy. Neither Theorem~\ref{thm:construct49} nor a faster representation proves a uniquely neural capability or a training-success probability for Adam or L-BFGS. The mathematical increment is the complete feasible, own-future construction and implementation account; the empirical question is its work and policy value relative to conventional generators.

\subsection{Original-owner compilation on a tensor cover}
On a tensor cover $\mathcal X=\prod_{j=1}^d\{x_{j,0},\ldots,x_{j,n_j}\}$, put $S=\prod_j(n_j+1)$. For arbitrary labels $y_i$, including inconsistent inexact labels, compute at each grid node $v$
\[
 (z_v,o_v)=\min_{\rm lex,i}(y_i+L\|v-x_i\|_1,i).
\]
The index $o_v$ is the original label owner, not the transformed node $v$.

\begin{theorem}[Continuous original-owner identity]\label{thm:compile49}
For a closed tensor cell $C$ with corners $\mathcal V(C)$, and $x\in C$,
\begin{align}\label{eq:compile49}
 f(x)&=\min_{v\in\mathcal V(C)}(z_v+L\|x-v\|_1),\\
 (f(x),i_*(x))&=\min_{\rm lex,v\in\mathcal V(C)}
 (y_{o_v}+L\|x-x_{o_v}\|_1,o_v).\notag
\end{align}
Coordinatewise forward and backward relaxations compute $(z,o)$ in $O(dS)$ additions and comparisons and $O(S)$ storage. After locating the cell, a query uses $O(d2^d)$ arithmetic operations. Exact rational comparisons preserve all ties, including $L=0$. Using the original owner's repaired action preserves the uncompiled policy, its cost and its certificate exactly.
\end{theorem}
\noindent The proof is in Supplement, Section~\ref{supp:proof-core49-3}.


The nodal algorithm is the classical distance transform \citep{felzenszwalb2012}; preserving the original action through continuous off-grid queries is the policy obligation. If there are $J_{t,i}$ node actions and $M_t$ integration representatives, write $Q_t=M_t\sum_iJ_{t,i}$. Literal dense cone evaluation costs $O(\sum_t dS_{t+1}Q_t)$ in representation work. Compilation replaces this by $O(\sum_t[dS_{t+1}+Q_td2^d])$, plus cell location, primitive evaluation, integration, feasible repair, certificate and output. Rational bit lengths and state boxes crossing many cells incur additional costs. There is no dimension-free complexity assertion.

\subsection{Acquired states and conventional actors}
A $b$-bit coordinate sensor reports a dyadic cell and its midpoint $\widehat x$. On the unit cube, $\|x-\widehat x\|_1\le r(b)=d2^{-b-1}$. For the investment capacity $\kappa(x)$ below, clip the original action to the smallest capacity throughout that cell and round down to an action quantum $\Delta_t$. This is feasible at every state in the cell. Relative to an exact true-state repair, its displacement is at most $2K_t^Ar(b)+\Delta_t$. Transporting the witness cone from $\widehat x$ to $x$ adds at most $2L_tr(b)$, plus any certified selector objective error $\nu_t$. Thus the implemented-policy account adds
\begin{equation}\label{eq:acquire49}
 \sum_{t<T}B_t\{(2L_t+2D_tK_t^A)r(b)+D_t\Delta_t+\nu_t\}.
\end{equation}
No continuity of the discrete owner selector is assumed. At uncertain numerical comparisons the simulator retains all possible actors, rather than choosing a favorable one.

For conventional FVI the continuation is multilinear on its realized tensor grid. Its actual modulus is bounded by the largest coordinate edge slope. A nearest-node actor uses a feasible node minimizer and true-state repair. Let $h_t$ be the largest half-cell $\ell^1$ radius, $L_t^{\rm gr}$ the Bellman graph modulus, and $e_t$ the node query allowance. Its optimal residual upper bound is $D_tk_t+e_t+L_t^{\rm gr}h_t$. Its selected-policy upper bound is $e_t+E_t$, where
\[
 E_t=\sup_x\{y_{i(x)}+L_t^{\rm gr}\|x-x_{i(x)}\|_1-f_t(x)\}.
\]
Partitioning cells by nearest-node boundaries makes the expression separately affine in each coordinate. Its maximum is therefore at a partition vertex; the implementation encloses all such vertices. Both uniform and curvature-refined FVI receive this sharper one-sided account, not an artificially doubled baseline. Their acquired-state allowance is derived separately in \eqref{eq:fvisensor49}.

\section{Original-optimum accuracy in the investment economy}\label{sec:accuracy62}
The exact witness theorem answers an implementation question conditional on a continuation. An economist also needs to know how close the complete implemented policy is to the original optimum. We now establish that connection in the investment economy itself. The argument exploits a property not used by the earlier cell-range calculation: despite the bilinear transition, the Bellman objective remains jointly convex. This permits a second-order interpolation account and transfers verified nodal decisions to an implementable policy over the entire state domain. It does not assume that a fitted neural continuation is convex.

\paragraph{Relation to the preceding two-control query.}
The earlier query experiment used the fixed shared capacity $1/8$. The complete dynamic extension below uses the original state-dependent capacity $c(x)$, with the retained exposure columns, action-cost interaction and second shock. Its nesting claim concerns the original scalar economic model; it is not an assertion that a full dynamic policy was already present in the earlier isolated query.

\subsection{Primitives and Bellman regularity}
Let $d\ge2$ be even, let indices be cyclic, and set $\beta=15/16$. Write
\begin{align*}
 S(x)&=\frac2d\sum_j(x_j-5/8)^2+
       \frac1{4d}\sum_j(x_j-x_{j+1})^2+
       2\left(\frac12-\frac2d\sum_jx_j\right)_+^2,\\
 g(x)&=S(x)+\frac2d\sum_j(x_j-5/8)^2,\qquad
 c(x)=\frac18+\frac1{8d}\sum_jx_j.
\end{align*}
The original scalar control has $0\le a\le c(x)$, running cost $S(x)+a^2+4a^4$, and transition
\[
 F_j(x,a,z)=\frac1{16}+\frac{x_j}{2}+\frac{x_{j+1}}8+
       \frac{x_j(1-x_{j+1})}{16}+B_ja+s_jz,
\]
where $(B_j,s_j)=(1/2,1)$ at alternating coordinates and $(1/4,-1)$ at the others, and $z$ is uniform on $[-1/32,1/32]$.

For the retained two-control extension, $a_1,a_2\ge0$, $a_1+a_2\le c(x)$, the second exposure column interchanges $1/2$ and $1/4$, and the action cost is
\[
 R(a)=a_1^2+a_2^2+4(a_1^4+a_2^4)+a_1a_2/4.
\]
An independent common-direction uniform innovation of radius $1/64$ is added. This is the previously specified extension, nested at $a_2=0$ when that innovation is disabled. Both models preserve $[0,1]^d$. Let $m\in\{1,2\}$ denote the number of controls and define
\begin{align*}
 Q_t^*(x,a)&=S(x)+R(a)+\beta\mathbb E V_{t+1}^*(F(x,a,Z)),\\
 V_t^*(x)&=\min_{a\in A(x)}Q_t^*(x,a),\qquad V_T^*=g.
\end{align*}
All costs and laws in this display are the economic primitives, not a fitted surrogate.

\begin{theorem}[Regularity of the original Bellman problem]\label{thm:regularity62}
For either control specification and every finite horizon, $V_t^*$ is convex and continuously differentiable in the interior, with the corresponding one-sided extension on the cube. It is $G_t$-Lipschitz in the $\ell^1$ norm and $M_t$-semiconcave in the Euclidean norm, where valid constants are
\begin{align*}
 G_T&=10/d,&G_t&=\frac{1037}{128d}+\beta\frac{47}{64}G_{t+1},\\
 M_T&=26/d,&M_t&=\frac{22}{d}+\frac{21}{256d}
             +\beta\left(\frac9{16}M_{t+1}+\frac18G_{t+1}\right).
\end{align*}
In particular, $G_t\le27/d$ and $M_t\le54/d$. The function $Q_t^*$ is jointly strongly convex on its convex feasible graph, with state curvature at least $107/(128d)$ and action curvature at least $7/4$.

A valid semiconcavity constant along each coordinate line is
\begin{align*}
 D_T&=\frac9d+\frac{16}{d^2},\\
 D_t&=\frac5d+\frac{16}{d^2}+\frac{21}{256d^2}
 +\beta M_{t+1}\left(\frac{85}{256}+\frac{23}{256d}\right).
\end{align*}
Consequently, for the positive multilinear interpolant $I_n$ on a mesh of spacing $1/n$,
\begin{equation}\label{eq:interpolation62}
 0\le I_nV_t^*(x)-V_t^*(x)\le\kappa_t,
 \qquad \kappa_t=\frac{dD_t}{8n^2}.
\end{equation}
\end{theorem}

\noindent The proof is in Supplement, Section~\ref{supp:proof-accuracy62-print-1}.


The theorem is a structural statement about these economic primitives. It is not a claim that every controlled diffusion, recursive preference specification, or game in the complete development edition has a convex Bellman problem. Nor does it require a signed-weight neural fit to preserve convexity. Convex approximation of dynamic programs has an established literature; the additional work here is to verify the needed curvature for the original nonlinear capacity-constrained model and connect it to the actual NBO actor and complete error account.

\subsection{Certified nodal Bellman bounds}
Suppose nodal arrays satisfy $l_{t+1}\le V_{t+1}^*\le u_{t+1}$ at all grid vertices. Equation~\eqref{eq:interpolation62} gives the functions $I_nl_{t+1}-\kappa_{t+1}$ and $I_nu_{t+1}$ as valid lower and upper continuations. The lower function need not itself be convex; convexity is used for the unknown true value and for the comparison argument, not asserted of arbitrary endpoint arrays.

Divide each uniform innovation component into $q$ equal bins and use its conditional mean. Convexity gives a lower quadrature bound for the true continuation. Semiconcavity gives an upper error
\begin{equation}\label{eq:quadrature62}
 \zeta_{t+1}=\frac{M_{t+1}}{6q^2}
                      \sum_{r=1}^{s}\rho_r^2\|s_r\|_2^2,
\end{equation}
where $s=1$ or $2$, $\rho_1=1/32$, and $\rho_2=1/64$ in the extension. The expectation remains the original continuous law. At the terminal date $g$ is evaluated directly, so no terminal interpolation allowance is needed.

At a vertex $v$, let $\mathcal A_A(v)$ be the capacity simplex lattice with spacing $c(v)/A$, including all its faces. For $m=1,2$, a minimizer can be represented by neighboring lattice vertices whose coordinate variances sum to at most $m c(v)^2/(4A^2)$. A valid action smoothness constant is
\[
 H_t=\frac{21}{4}+\beta M_{t+1}\|B\|_2^2,
 \quad \|B\|_2^2=\frac{5d}{32}\ (m=1),\quad
 \|B\|_2^2=\frac{9d}{32}\ (m=2).
\]
Taylor's inequality and that barycentric representation imply
\begin{equation}\label{eq:action-cover62}
 0\le\min_{a\in\mathcal A_A(v)}Q_t^*(v,a)-V_t^*(v)
       \le\delta_t:=\frac{mH_t}{128A^2}.
\end{equation}
The linear Taylor terms cancel even when the optimum lies on a face. This is an error against the continuous feasible optimum, not only against the finite action menu.

Let $q_t^-(v,a)$ and $q_t^+(v,a)$ include all interpolation, quadrature and arithmetic errors and enclose $Q_t^*(v,a)$. Then
\begin{equation}\label{eq:nodal62}
 l_t(v)=\max\{0,\min_{a\in\mathcal A_A(v)}q_t^-(v,a)-\delta_t\},
 \qquad
 u_t(v)=\min_{a\in\mathcal A_A(v)}q_t^+(v,a)
\end{equation}
are valid nodal Bellman bounds. The zero truncation uses the nonnegative economic costs. Extra neural or conventional proposals may improve the upper side but do not remove the lower action-cover allowance. Positive interpolation, averaging and minimum maps are monotone and nonexpansive; the separate allowances in \eqref{eq:interpolation62}--\eqref{eq:action-cover62} are therefore propagated, rather than being discarded because an endpoint calculation contains the truth.

\subsection{From a fitted witness to the complete implemented policy}
Let $a_t(v)\in A(v)$ be the actual action returned at each vertex by any generator. It may be selected by the exact neural witness, a conventional fitted continuation, or their common verified fallback. Write $w_v(x)$ for the multilinear weights and set $\bar\pi_t(x)=\sum_vw_v(x)a_t(v)$. Affineness of capacity makes this action feasible. Let $q_t^{\mathrm{sel},+}(v)$ be a verified upper bound on its original-optimum continuation value $Q_t^*(v,a_t(v))$, using the same lower and upper Bellman reference for all generators.

\begin{theorem}[Actual-policy, whole-domain certificate]\label{thm:policy62}
Set
\[
 e_t=\max_v\{q_t^{\mathrm{sel},+}(v)-l_t(v)\}+\kappa_t.
\]
Suppose the implemented policy $\widehat\pi_t$ is feasible and differs from $\bar\pi_t$ by at most $b_t$ in action $\ell^1$ norm, uniformly over the state cube. Define
\[
 \Lambda_t=\frac{13}{16}+\beta\frac{3d}{8}G_{t+1},\qquad
 D_T^\pi=0,\quad D_t^\pi=e_t+\Lambda_tb_t+\beta D_{t+1}^\pi.
\]
Then, at every true state and date,
\[
 0\le J_t^{\widehat\pi}(x)-V_t^*(x)\le D_t^\pi.
\]
This statement is conditional on the actual recorded policy, not on success of its fitting algorithm. It applies equally to a pure fitted policy and to a policy augmented by verified common candidates.
\end{theorem}

\noindent The proof is in Supplement, Section~\ref{supp:proof-accuracy62-print-2}.


In the execution, observation is rounded down at forty bits. An acquired state $y\le x$ has no greater capacity than the true state. Each interpolated action is then rounded down at twenty bits using the lower endpoint of its arithmetic enclosure. If $L_{ij}$ bounds the derivative of action coordinate $i$ with respect to state coordinate $j$, and $\eta$ bounds native interpolation error, a valid implementation allowance is
\[
 b_t\le 2^{-40}\sum_{i,j}L_{ij}
           +m(2^{-20}+2\eta+\eta_{\mathrm{out}}).
\]
Here $L_{ij}$ is computed from the maximum adjacent nodal difference times $n$, and $\eta_{\mathrm{out}}$ covers endpoint subtraction. All these numbers are recorded outward. This new barycentric acquired actor is explicitly distinguished from the earlier discontinuous cell-constant actor; an old actor is not silently replaced during a purported fixed-policy revalidation.

\subsection{Arithmetic, work, and scope}
The native interpolation kernel forms nonnegative tensor weights and a weighted sum of at most $2^d$ stored binary64 endpoints. It is compiled without reassociation or floating-point contraction. For $K=(3d+4)2^d+32$, an explicit conservative absolute allowance is
\[
 \eta=\frac{K2^{-53}}{1-K2^{-53}}\max_v|z_v|+K2^{-1022}.
\]
The final term conservatively covers underflow. Input coordinates, affine indexing, finite endpoints and the binary64 format are checked. Roundoff expansion and the ordinary interval cost arithmetic remain separate from the analytic approximation bounds. Exact rational fixtures compare the native sum with its mathematical tensor interpolant.

With $N=(n+1)^d$ vertices, $K_A=A+1$ scalar actions or $K_A=(A+1)(A+2)/2$ two-control actions, the reference calculation uses $O(TNK_Aq^s2^d d)$ elementary work under the explicit tensor implementation. Stored reference values use $O(TN)$ words and a policy uses $O(TNm)$ words. Streaming batches avoid a state-pair table. Proposal fitting, exact rational witness work, actor transfer, serialization, independent cost inference and compilation are additional costs, not absorbed into this bound. The analytic errors are second order in state, action and innovation resolution; the state-count and tensor-evaluation factors still depend exponentially on dimension. The result supplies a matched original-optimum comparison without asserting that these costs have disappeared.

\section{Neural Bellman queries without a state lattice}\label{sec:query63}
The preceding construction makes numerical error accountable, but a complete tensor reference can absorb the benefit of a fitted continuation. We therefore change the computational organization, not the economic problem. A Neural Bellman selector proposes an action at the state actually acquired. A recursive interval query then bounds the original Bellman value at that state and checks the proposed action against the entire continuous feasible interval. No conventional policy is precomputed at all states. The controller is the resulting callable algorithm, not the unverified proposal alone.

The distinction matters for attribution. Learning may reduce the work needed to locate a useful action. It does not supply the inequalities used to certify that action. The adaptive conventional comparator receives the same inequalities and the same original-optimum target. The trained value-fit policies in the preceding study remain distinct from the action-query predictor introduced here.

\subsection{A curvature lower bound on an adaptive action partition}
Use the original scalar economy of Section~\ref{sec:accuracy62}, with $r$ dates remaining. Write $V_r$ for its optimal value and $Q_r(x,a)$ for the value of a first action followed by the optimal future. Let $G_r,M_r$ be the primitive constants in Theorem~\ref{thm:regularity62}, indexed by remaining horizon, and put
\[
 H_r=\frac{21}{4}+\beta M_{r-1}\frac{5d}{32},
 \qquad \Lambda_r=\frac{13}{16}+\beta G_{r-1}\frac{3d}{8}.
\]
Thus $H_r$ is an upper curvature bound in the scalar action, and $\Lambda_r$ bounds the absolute action derivative. Neither depends on a fitted model.

\begin{lemma}[Parabolic minorants]\label{lem:parabolic63}
Suppose $f$ is $H$-semiconcave on $[l,u]$, and $L_l\le f(l)$, $L_u\le f(u)$. For $s\in[0,1]$,
\[
 f(l+s(u-l))\ge (1-s)L_l+sL_u
                 -\frac H2(u-l)^2s(1-s)=:p(s).
\]
Consequently $\min_{[0,1]}p$ is a lower bound for the minimum of $f$ on the interval. With $K=H(u-l)^2/2$ and $D=L_u-L_l$, this lower bound is $L_l$ if $D\ge K$, $L_u$ if $D\le-K$, and otherwise
\[
 \frac{L_l+L_u}{2}-\frac K4-\frac{D^2}{4K}.
\]
The endpoints $L_l,L_u$ need not themselves belong to a convex interpolant. Increasing $K$ preserves the lower-bound property.
\end{lemma}
\noindent The proof is in Supplement, Section~\ref{supp:proof-query63-1}.


At an acquired state, retain an ordered partition of the feasible action interval and an outward value bracket at every partition point. The minimum of the parabolic minorants bounds the continuous-action infimum from below. The smallest upper endpoint supplies an actual feasible witness. A neural or polynomial proposal inserts another point into this partition; it cannot remove the endpoints or any part of the feasible interval. The running lower bound is the maximum of all previously valid global lower bounds, so new information cannot weaken the certificate. The upper bound is the minimum over every evaluated feasible witness.

\subsection{Recursive queries under the original continuous law}
For an endpoint action with $r>1$, divide the scalar innovation into $q_r$ equiprobable bins and query $V_{r-1}$ at each conditional-mean next state. If each child returns width at most $\eta_r$, the endpoint bracket has width at most
\begin{equation}\label{eq:querywidth63}
 \beta\eta_r+\frac{\beta M_{r-1}d}{6144q_r^2}+\omega_r.
\end{equation}
Here $\omega_r$ includes outward primitive arithmetic and the Lipschitz expansion from an interval-valued next state to its queried center. Convexity provides the lower conditional-mean inequality; semiconcavity and the exact conditional variance provide the upper allowance. The innovation remains continuous. There is no state-interpolation term because a child value is queried rather than read from a state lattice.

For a requested local gap $\delta$, the implementation sets $\eta_r\le\delta/(4\beta)$ and chooses a dyadic $q_r$ so that the second term in \eqref{eq:querywidth63} is at most $\delta/4$. A separate endpoint-width guard charges arithmetic and the small capacity-rounding remainder. Failure of that guard is a precision failure, not a successful return. The terminal query uses analytic integration instead of recursive quadrature.

\begin{theorem}[A state-grid-free NBO service]\label{thm:query63}
Consider the original scalar investment economy. At every recursive query, suppose the outward endpoint and capacity-transfer accounts give width at most $5\delta/8$. Let the lower bound be the running maximum described above. Refine an interval attaining the smallest current minorant whenever the global upper-minus-lower gap exceeds $\delta$. Each new point must lie between one fifth and four fifths of that interval. Arbitrary additional learned proposals may be inserted initially. Then:

(i) every returned interval $[L_r(x),U_r(x)]$ contains $V_r(x)$, and its returned feasible action $a_r(x)$ satisfies
\[
 Q_r(x,a_r(x))-V_r(x)\le U_r(x)-L_r(x)\le\delta;
\]
(ii) the real-arithmetic refinement terminates. With at most one interior learned proposal, a sufficient point budget is
\begin{equation}\label{eq:querycap63}
 6+\left\lceil5c(x)\sqrt{H_r/(3\delta)}\right\rceil;
\end{equation}
(iii) for the policy that calls this service at every subsequent date, the original-law loss satisfies
\begin{equation}\label{eq:querypolicy63}
 J_T^{\pi}(x)-V_T(x)
 \le\sum_{t=0}^{T-1}\beta^t
 \left[\delta_t+
 \left(8+\frac{11\beta dG_{T-t-1}}{16}
                  +dG_{T-t}\right)2^{-b}\right],
\end{equation}
when the state is acquired downward at $b$ bits per coordinate and every returned action is feasible at that acquired state. All action rounding must precede its endpoint evaluation. The inequalities hold uniformly in the true initial state, not only at the states used in an experiment.
\end{theorem}
\noindent The proof is in Supplement, Section~\ref{supp:proof-query63-2}.


The finite-precision contract is explicit. The executed service uses the width guard, a 512-point cap, and a 32-bit downward action lattice. A precision or cap failure withholds the floating certificate and invokes a separately counted exact-rational enumeration. This fallback constructs values only at the required descendants; it does not load a precomputed conventional policy. The rational fallback chooses dyadic action and quadrature counts until each cover or quadrature allowance is at most one quarter of the request, and requests child width at most one quarter divided by the discount. Provided $\delta>8\Lambda_r2^{-32}$ at every recursive level, its exact brackets have width at most three quarters of the request. This proves a finite return for every acquired state in the declared tasks, including states absent from the workload. Rational bit-operation cost is additional work and is not replaced by a unit-cost floating clock. The mathematical point budget is checked against the declared tolerance range; termination in an arbitrary floating implementation is not inferred merely from the real-arithmetic argument. Fitting convergence is unnecessary for validity: any finite prediction, and even its clipped replacement when numerically invalid, is only a proposed query.

\subsection{A certified analytic terminal kernel}
The last decision admits an especially strong common comparator. Let $y=F(x,0,0)$, $b_0=1/2-2\operatorname{mean}(y)$, and $\gamma_j=B_j$. Define
\[
 D(x)=\frac8d\sum_j\gamma_j(y_j-5/8)
   +\frac1{2d}\sum_j(\gamma_j-\gamma_{j+1})(y_j-y_{j+1}).
\]
For every even $d$, the exact terminal action derivative is
\begin{equation}\label{eq:cubic63}
 16a^3+\left(2+\frac{41\beta}{32}\right)a
 +\beta D(x)-3\beta(b_0-3a/4)_+.
\end{equation}
Each branch is a strictly increasing cubic. Its floating closed-form root supplies only a candidate. The independently evaluated outward derivative and exact-law objective certify that candidate; a safeguarded derivative bracket remains available if the candidate is inaccurate.

For completeness, if $f$ is $\mu$-strongly convex and an outward interval encloses $f'(a)$, the largest compatible projected derivative residual $R$ gives
\[
 f(a)-\min_{[0,c]}f\le R^2/(2\mu).
\]
At zero, a nonnegative derivative has residual zero; at capacity, a nonpositive derivative has residual zero. Otherwise the absolute derivative bound is sufficient. This follows by minimizing the strong-convexity supporting quadratic over the feasible interval and, conservatively, over the line. Here $\mu=2$. Equation~\eqref{eq:cubic63}, its $41/32$ and $9/4$ coefficients, analytic shock moments, and outward derivatives are independently checked against rational formulas. Every method uses this kernel; it is not credited to the neural predictor.

\subsection{Work, prediction, and the economic object}
Let $A_r(\delta)$ be a bound on action endpoints at a query. Excluding training, a scalar operation account obeys
\[
 W_1(\delta)=O(dA_1),\qquad
 W_r(\delta)\le C d A_r(\delta)+
 A_r(\delta)q_r W_{r-1}(\delta/(4\beta)).
\]
The constants $dG_r,dM_r$ are bounded independently of $d$, so the declared endpoint and quadrature budgets need not grow exponentially with the state dimension. Primitive evaluations still cost $O(d)$. This removes the state lattice, not the horizon or shock-tree cost. The executed vectorization batches a query's innovation descendants. Its largest live state batch is reported; it does not claim the smaller memory of a fully streamed implementation. Increasing the horizon or the shock dimension may still be expensive. The new execution is scalar-control; the earlier full two-control result is retained rather than relabeled as a query-service experiment.

A useful prediction can reduce the endpoint count or return its own action from the initial partition. An inaccurate prediction can instead add work. This is a prediction-dependent algorithm, not a theorem of neural dominance. Learned warm starts have their own literature, including \citet{sakaue2023warm}; their discrete-convex setting differs from the original-law Bellman intervals here. Convex control policies and learned optimization parameters are studied by \citet{agrawal2020control}. Our comparison identifies the extra proposal cost and the resulting query transcript rather than claiming priority for learning a warm start.

The purchase is now a complete callable controller with a declared error contract. Its training cost, every subsequent certified query, and its return-time cost belong to the same service. A reduction in queried state storage is economically relevant for repeated decisions only after these costs are counted. The original cost normalization is maintained. Workload seconds and normalized expected-cost loss are reported separately; a price per unit of computation can convert the former to the latter only as an explicitly specified decision-maker parameter.

\section{Certificate-gated neural queries}\label{sec:gated64}
A prediction of the optimal action and a prediction of a useful numerical query are different objects. Inserting a learned action before examining an already available certificate charges an extra continuation query even when the two boundary actions suffice. We instead make the demand for another query a mathematical decision and let a trained NBO selector determine its location only after that demand has been established. The underlying continuation recursion and the economic optimum are unchanged.

At a recursive state, start with the two feasible action endpoints. Compute the global lower bound from the parabolic minorants of Lemma~\ref{lem:parabolic63}, including the capacity remainder, and the smallest upper endpoint. Return immediately when their difference is at most the requested tolerance. Otherwise choose a partition interval with smallest current minorant. A learned action is used only if it lies in the central half of that interval; it then replaces the conventional proposed split. Outside that region the safeguarded parabolic split is used. The rounded split must pass the one-fifth interior check before evaluation. A prediction is evaluated at most once at each recursive state and is reused only as a proposed location. It supplies neither a value endpoint nor a lower bound.

\begin{theorem}[Gated prediction and complete-policy validity]\label{thm:gated64}
Under the assumptions of Theorem~\ref{thm:query63}, the gated service returns the same original-optimum and actual-policy guarantees for every supplied prediction, independently of fitting error. In real arithmetic, with endpoint-plus-capacity width at most $5\delta/8$, a sufficient number of endpoint evaluations at a state is
\[
 A_r(\delta,x)=4+\left\lceil5c(x)\sqrt{H_r/(3\delta)}\right\rceil.
\]
Each requested evaluation either reduces the unresolved action partition or terminates the query; no initial interior evaluation is mandatory. If the boundary certificate suffices, no predictor is evaluated at that state. If every prediction is declined at every recursive query, the value endpoints and implemented actions coincide with those of the conventional parabolic service using the same arithmetic and tie rules.
\end{theorem}
\noindent The proof is in Supplement, Section~\ref{supp:proof-gated64-1}.


This theorem establishes robustness to the numerical content of a fitted prediction, not that learning is free or always useful. It also does not state that an arbitrary routed transcript has fewer queries than the adaptive conventional transcript. The executed comparison records both possibilities. Unlike a policy catalogue with a stored reference action at every state, neither service constructs a conventional state lattice before it can return a certified policy.

\subsection{State dimension, working memory, and reliability}
For the original scalar model, the primitive bounds $dG_r\le27$ and $dM_r\le54$ imply
\[
 H_r\le\frac{3369}{256},\qquad
 \Lambda_r\le\frac{1319}{128},\qquad
 \frac{\beta M_{r-1}d}{6144}\le\frac{135}{16384}.
\]
Thus the action and scalar-innovation budgets at fixed remaining horizon and tolerance are independent of the number of state coordinates. Every primitive evaluation still requires $O(d)$ arithmetic. A width-$w$ one-hidden-layer ReLU prediction on the recorded $d+5$ features requires $O(dw)$ work and stored parameters; a full quadratic prediction on these features uses $O(d^2)$ terms. Neither cost is included for free in the Bellman recursion.

Let $q_r(\delta)$ be the dyadic quadrature count and let $P_r(d)$ be the cost of one predictor evaluation. Apart from the separately identified exact-arithmetic recovery, a conservative work recursion is
\[
 W_r(d,\delta)\le C d A_r(\delta)+P_r(d)
   +A_r(\delta)q_r(\delta)
      W_{r-1}(d,\delta/(4\beta)).
\]
The predictor term is absent on a boundary-certified query. This expression removes the exponential state lattice, not the horizon or innovation-tree dependence. A depth-first scalar implementation can retain the current action partition and one innovation child at a time. The executed batched implementation instead reports its maximum state batch and process peak memory. A zero count of stored state-lattice nodes does not mean zero storage.

The deterministic certificate is valid for any supplied finite network weights. A failure to fit and return such weights is a separate service event. Our reliability estimand is the return count and worst measured complete work over the preregistered finite seed catalogue, with every repetition retained. Numerical workload paths are not optimizer replications. The real-arithmetic cap and the exact-rational precision condition establish algorithmic termination at admissible tolerances; recorded processor seconds are descriptive, not a hardware-uniform worst-case bound.

\subsection{An explicit economic purchase criterion}\label{sec:purchase64}
The decision maker buys a callable controller rather than a table of successful sample decisions. Fix a physical capital numeraire for the normalized state, a conversion $m>0$ dollars per unit of the model's loss index, a computation charge $p\ge0$ dollars per measured second, and a fixed deployment fee $f\ge0$ dollars. For $n$ uses under a specified initial-state law, the monetary objective is
\[
 \mathcal C(\pi)=nm\,\E[J_0^\pi(X_0)]+pW(\pi)+f.
\]
These are declared preference and service-price parameters, not estimated economic primitives. The policy theorem bounds the expected loss under the original continuous law for every such initial-state law. The finite dyadic workload below measures execution; its average cost is not substituted for that expectation.

If two policies have an independently valid expected-gain interval $[l,u]$ for $\E[J_0^{\pi_0}-J_0^{\pi_1}]$, and their incremental measured service charge is $\Delta=p(W_1-W_0)+f_1-f_0$, then the monetary net gain lies in
\[
 [nm l-\Delta,\;nm u-\Delta].
\]
This follows by multiplication by the positive numeraire conversion and subtraction of the same service charge. A positive lower endpoint supports purchase; a negative upper endpoint rules it out at the stated prices. Uniform regret certificates alone do not supply a directional expected-gain interval. They instead support a worst-case accuracy requirement imposed before selecting the least expensive eligible service.

Finally, a forecast of amortized query work is distinguishable from a measured complete service. With a fixed model and exactly repeated query workload, if training requires $B$ operations and deployment requires $D_L$ rather than $D_C$ operations per workload, then learning uses fewer total operations after $k$ repetitions precisely when $D_L<D_C$ and $k>B/(D_C-D_L)$. If $D_L\ge D_C$ and $B\ge0$, no repetition of that unchanged workload produces an operation saving. This arithmetic criterion includes training and cannot be applied to unrelated future states without a new assumption. The executed catalogue reports its own full process clocks, not a favorable hypothetical amortization.

\section{When a predicted query completes the certificate}\label{sec:localization65}
The safeguarded theorem is insensitive to the quality of a supplied prediction. A separate question is when that prediction is sufficient to finish verification. The answer must concern both the selected action and the unresolved lower bound. Proximity to a minimizer alone does not make a coarse lower certificate accurate.

\begin{proposition}[Witness accuracy and certificate localization]\label{prop:localization65}
Let $f$ be a differentiable, $H$-semiconcave objective with $H\ge0$ on $[0,c]$, let $a^*$ minimize $f$, and suppose $|f'|\le\Lambda$. An evaluated proposal $\widehat a$ lies in $[0,\underline c]$, where $0<\underline c\le c$ and $c-\underline c\le\Delta$. Each endpoint interval is a valid enclosure of $f$ and has width at most $w$. Suppose the partition containing $\widehat a$ covers $[0,\underline c]$ and has maximum interval length $h$. Let $L$ be its parabolic lower certificate with the capacity allowance $\Lambda\Delta$, and let $U$ be the smallest upper endpoint, retaining earlier valid bounds as in the query algorithm. With $e=|\widehat a-a^*|$ and $\lambda=|f'(a^*)|$,
\begin{equation}\label{eq:localization65}
 U-L\le 2w+\Lambda\Delta+\frac{Hh^2}{8}
                         +\lambda e+\frac{He^2}{2}.
\end{equation}
For an interior minimizer $\lambda=0$. At a boundary the linear term cannot generally be removed. Thus the right side at most $\delta$ is a sufficient certificate-return condition, not an assumption that a trained prediction meets it.

If old and new valid endpoint records give $(L,U)$ and $(L',U')$ after one additional query, with the running maximum lower bound and minimum upper bound, then
\begin{equation}\label{eq:certificate-progress65}
 (U-L)-(U'-L')=(U-U')+(L'-L)\ge0.
\end{equation}
An added prediction may have zero certified progress. The same identities apply to a conventional query.
\end{proposition}
\noindent The proof is in Supplement, Section~\ref{supp:proof-localization65-1}.


This proposition isolates what a fitted representation can contribute without crediting it with a conventional solution. It can provide a small-error witness and change which unresolved interval is refined. The global partition, endpoint arithmetic and capacity account still have to support that witness. The maximum-width sufficient condition is deliberately conservative; the implemented minimum-minorant rule can stop earlier when only a subset of the partition remains competitive.

\subsection{A complete accuracy contract and its resource dependence}
For $T$ dates and target $\epsilon$, write $S_T=\sum_{t<T}\beta^t$. The recorded service chooses a common local tolerance by rounding $\epsilon/(2S_T)$ to binary64 and verifies the final outward allowance. Under $dG_r\le27$, a uniform per-date acquisition allowance at forty bits is
\[
 \alpha_b=\left(8+\frac{297\beta}{16}+27\right)2^{-40}.
\]
The service certifies $(\delta+\alpha_b)S_T\le\epsilon$. Each descendant request is $\delta/(4\beta)$, rounded as in the actual code and included in its outward endpoint guard. At the smallest request in the two- and three-date catalogues, the declared thirty-two-bit action quantum satisfies the exact-recovery condition $\delta_r>8\Lambda_r2^{-32}$. The supplement checks this inequality with rationals for every horizon and target, rather than inferring it from a successful workload.

For fixed remaining horizon and accuracy, the bounds on action endpoints and scalar-innovation bins are independent of state dimension in this normalized economy. The factors $dG_r$ and $dM_r$ are bounded, while every transition and cost still reads $d$ coordinates. A width-$w_\theta$ one-layer ReLU predictor has $O(dw_\theta)$ parameters and evaluation work; a dense quadratic on the same features has $O(d^2)$ terms. These statements identify the absence of a $(n+1)^d$ stored state table. They do not count an arbitrary-precision rational operation as a unit-cost floating operation and do not suppress the product of descendant query counts as horizon grows.

The two fitted implementations and their conventional controls therefore share one original-optimum contract, but need not share a transcript, a returned policy or a runtime. The completed experiments charge their distinct transcripts. A finite-catalogue reliability statement concerns all specified seeds and repetitions, including fitting warnings. It is not a population statement about optimizer success.

\section{Operational prediction certificates and readout training}\label{sec:operational67}
Proposition~\ref{prop:localization65} separates witness accuracy from lower-certificate localization. Its distance to an unknown minimizer is useful for analysis but is not a stopping observable. We now replace that distance by an upper envelope computed from information already purchased. This gives a certificate before an additional Bellman query, a training criterion measured in unresolved verification work, and a constructive neural readout with a computable optimization gap.

\subsection{A decision before the next query}
Fix a state, a date, and the current action partition. Write $f(a)=Q_t^*(x,a)$, and suppose $f$ is $\mu$-strongly convex on the feasible interval. Let $\overline f_l,\overline f_u$ be verified upper endpoints at two consecutive actions $l<u$. Let $L$ be the current lower bound for the original feasible infimum, including the capacity and arithmetic allowances. For $a=l+s(u-l)$, define
\begin{equation}\label{eq:screen67}
 C(a)=(1-s)\overline f_l+s\overline f_u
       -\frac\mu2 (a-l)(u-a),\qquad 0\le s\le1.
\end{equation}
Every symbol in this expression is either a primitive bound or an existing endpoint. No derivative or action at the unknown optimum occurs.

\begin{theorem}[Pre-query closure]\label{thm:screen67}
Under the preceding conditions, $f(a)\le C(a)$ throughout $[l,u]$. Let $a_\Delta\in[l,u]$ be the actual feasible rounded proposal. An outward upper endpoint $\overline C(a_\Delta)$ satisfying
\begin{equation}\label{eq:screen-stop67}
 \overline C(a_\Delta)-L\le\delta
\end{equation}
certifies $f(a_\Delta)-\inf_{a\in A(x)}f(a)\le\delta$ without another Bellman query. If the screen does not close, the protected refinement of Theorem~\ref{thm:gated64} retains its validity and finite query bound. The complete acquired policy therefore retains the original-optimum account of Theorem~\ref{thm:query63}.

Equivalently, the union over partition intervals of the sublevel sets $\{a:C(a)-L\le\delta\}$ is a computable certified action set. Each constituent is a quadratic sublevel interval intersected with its partition cell. Endpoints obtained numerically must be rounded inward, or the final action must be checked by \eqref{eq:screen-stop67}.
\end{theorem}
\noindent The proof is reproduced in Supplement, Section~\ref{supp:relocated69-1}.

The screen is shared by every comparator. In particular, minimizing the quadratic upper envelope is an available conventional proposal. Prediction is valuable only to the extent that its proposed location or evaluation cost improves the observed service. The executed neural rule screens its feasible prediction on the partition interval containing it, even when that location is outside the protected part of the currently critical interval. A failed screen does not authorize an unprotected split.

\subsection{A computable loss for unresolved work}
A \emph{certificate context} $z$ consists of the state, remaining horizon, tolerance, partition endpoints and their verified bounds. A proposal $a_\theta(z)$ is feasible and is computed without another Bellman query. Define its verified score and a margin loss by
\begin{align}
 s_\theta(z)&=\overline C_z(a_\theta(z))-L_z,\label{eq:score67}\\
 \ell_\tau(z,\theta)&=\min\left\{1,
       \frac{(s_\theta(z)-\delta_z+\tau)_+^2}{\tau^2}\right\},
       \qquad \tau>0.\label{eq:work-loss67}
\end{align}
Let $K(z)$ be a valid cap on further endpoint queries if the screen fails and the protected fallback is followed. This is a count of endpoint queries at the current state. The recursive cost of each endpoint is charged through the earlier descendant-work recursion. On contexts satisfying the declared precision contract, the real-arithmetic and rational-recovery bounds supply finite such caps.

\begin{theorem}[Certificate loss and additional work]\label{thm:work-loss67}
If $N_\theta(z)$ is the number of additional endpoint queries after the screen, then
\begin{equation}\label{eq:work-bound67}
 N_\theta(z)\le K(z)\mathbf1\{s_\theta(z)>\delta_z\}
                    \le K(z)\ell_\tau(z,\theta).
\end{equation}
Consequently, an achieved margin loss bounds the average additional queries on any specified finite context catalogue, with no unknown minimizing action.

For a distributional assertion, condition on construction of a finite family of $M$ predictors. Suppose $z_1,\ldots,z_n$ are independent validation contexts from a declared predictor-independent law, and $0\le K(z)\le K_{\max}$. With probability at least $1-\alpha$, simultaneously for that family,
\begin{equation}\label{eq:work-validation67}
 \mathbb E N_\theta(z)\le
 \frac1n\sum_{i=1}^n K(z_i)\ell_\tau(z_i,\theta)
 +K_{\max}\sqrt{\frac{\log(M/\alpha)}{2n}}.
\end{equation}
The cost of obtaining and validating contexts and of evaluating the predictor is additional work, not part of $N_\theta$.
\end{theorem}
\noindent The proof is reproduced in Supplement, Section~\ref{supp:relocated69-2}.

The context-law requirement matters. A stream of recursively generated contexts can depend on earlier predictions; it is not automatically an independent validation sample for \eqref{eq:work-validation67}. The deterministic inequality \eqref{eq:work-bound67} remains valid context by context. The recorded action-label mean squared error is not substituted for \eqref{eq:work-loss67}. The new loss measures a certificate's unresolved work rather than resemblance to an approximate action label.

\subsection{A constructive ReLU readout with an optimization certificate}
The work loss can be trained without invoking a successful hidden-layer optimizer. Fix a nonnegative affine--ReLU feature map $\phi(z)\in\mathbb R^p$ with $\sum_j\phi_j(z)=1$. Linear hat functions on a fixed coordinate partition are one exact construction; averages of such maps give multivariate features. A box-constrained readout $w\in[0,1]^p$ proposes
\[
 a_w(z)=l_z+(u_z-l_z)\phi(z)^\top w.
\]
The partition interval and features are fixed before optimizing $w$. Let $\eta_i\ge0$ cover action rounding and evaluation of the score, so that the computed score of the actual rounded action is at most $C_i(a_w(z_i))+\eta_i-L_i$. For example, on lattice-aligned endpoints an action Lipschitz bound gives $\eta_i=\Lambda_i\Delta_a$ for downward rounding. Define
\begin{equation}\label{eq:readout-loss67}
 F(w)=\frac1n\sum_i K_i
 \bigl(C_i(a_w(z_i))+\eta_i-L_i-\delta_i+\tau\bigr)_+^2
                      +\frac\lambda2\|w\|_2^2,
 \qquad \lambda\ge0.
\end{equation}
Unlike an action-label objective, this loss directly upper-bounds the additional-work account. It trains the upper witness against the current global lower certificate.

\begin{proposition}[Certified readout training]\label{prop:readout67}
The function $F$ is convex and continuously differentiable on the box. For $g=\nabla F(w)$, put
\begin{equation}\label{eq:dual-gap67}
 D(w)=g^\top w-\sum_{j=1}^p\min\{0,g_j\}.
\end{equation}
Then $0\le F(w)-\min_{v\in[0,1]^p}F(v)\le D(w)$. On the training contexts the average additional endpoint-query count is at most $F(w)/\tau^2$. If $\lambda>0$, the minimizer $w^*$ is unique and
\[
 \|w-w^*\|_2\le\sqrt{2D(w)/\lambda}.
\]
All quantities in these conclusions can be bounded by outward arithmetic; they are exact rational quantities when the contexts, features and iterate are rational.
\end{proposition}
\noindent The proof is reproduced in Supplement, Section~\ref{supp:relocated69-3}.

This is a fixed-feature neural construction, not a convergence assertion for unrestricted nonconvex network fitting. The deposited rational implementation uses finite conditional-gradient steps and reports its achieved loss and dual gap. Its regression examples test the proposition; they are not new investment-policy observations. The service comparison below retains the preregistered L-BFGS and conventional fits, rather than relabeling those models as minimizers of \eqref{eq:readout-loss67}. Prediction-dependent optimization has an established literature \citep{sakaue2023warm,sambharya2023}; the additional object here is a computable original-law Bellman certificate that can license an action before another query and supplies a directly auditable training criterion.

\section{Training, validation, and complete work}\label{sec:training69}
The fixed-feature proposition is now executed as its own numerical arm, rather than being attributed to the earlier hidden-layer fits. For $x\in[0,1]^d$, compare the two nonnegative feature maps
\[
 \phi_{jh}^{\rm R}(x)=d^{-1}(1-|4x_h-j|)_+,
 \qquad
 \phi_{jh}^{\rm B}(x)=d^{-1}{4\choose j}x_h^j(1-x_h)^{4-j},
 \quad j=0,\ldots,4.
\]
Both sum to one and have $5d$ coefficients. The first is an affine--ReLU hat construction; its circuit and its hat basis are one method. The second is the conventional degree-four Bernstein construction. Both readouts receive the same purchased states, boundary Bellman enclosures and approximate labels, the same certificate-aware objective, box constraint, penalty, iteration cap and weight lattice. A label-loss ReLU control retains the same features but changes the objective. The quadratic control receives the same information, with its different coefficient count reported. Prediction-free controls incur no unused training cost.

After quantizing the returned weights, we recompute $F(w)$ and the box dual gap $D(w)$ exactly from rational contexts and features. This certifies the achieved iterate even if the finite optimizer reports nonconvergence. The rounding allowance is checked against actual rounded training proposals. Future policy returns still require the outward screen; neither a small training loss nor the optimizer's status can authorize an action.

Validation uses a disjoint, predictor-independent stream of uniform acquired-state contexts. A failed first screen invokes a declared complete-action mesh with a finite root-query cap $K(z)$. This recovery and its descendants are executed. Its context law is not the distribution of states visited by recursive deployment, so the distributional validation conclusion is not silently transferred to those states.

\begin{proposition}[A finite catalogue of complete work]\label{prop:complete-work69}
Condition on a frozen finite family of predictors and a predictor-independent validation law. Suppose a failed screen invokes a deterministic recovery whose root-query count is at most $K(z)$, and every root query costs at most $W(z)$ under the declared descendant and arithmetic contract. If $0\le W(z)K(z)\le C_{\max}$, then with probability at least $1-\alpha$, simultaneously for $M$ fixed predictors,
\[
 \mathbb E[\text{recovery work}_\theta]
 \le\frac1n\sum_{i=1}^n W(z_i)K(z_i)\ell_\tau(z_i,\theta)
       +C_{\max}\sqrt{\frac{\log(M/\alpha)}{2n}}.
\]
Training, context acquisition, prediction, loading, identity checks and output costs are additional terms in complete work. Replacing $W$ by one gives a root-query result, not a bound on elapsed process time.
\end{proposition}
\begin{proof}
On a failed screen the clipped loss equals one and recovery work is at most $WK$; on success recovery work is zero. Apply the bounded-sum inequality to $WK\ell_\tau\in[0,C_{\max}]$ and take a union bound. The separately named service costs are not recovery work.
\end{proof}

The executed validation uses the root-query version. Complete process clocks remain separate measurements. Each reuse process fits and serializes once, then executes sixty-four declared blocks. Model loading checks the source identity, task, horizon, acquisition precision and stored hash. Block sixteen deliberately presents a stale acquisition identifier and uses conventional verification. An early persistent crossing requires cumulative measured work below the matched adaptive process by block forty-eight and at every later block through sixty-three, using zero-based indices. A last-block crossing does not qualify. Four process repetitions describe timing variability; they do not establish an infinite-horizon crossing or a hardware-invariant probability.

For two controls, a trained two-output proposal supplies a paid feasible upper witness while the original triangular lower cover remains intact. Projection and downward rounding precede its original-law evaluation.
\begin{proposition}[Optional vector witnesses]\label{prop:vector-witness69}
Suppose a triangular action cover has a valid global lower bound $L$ and feasible evaluated witnesses with upper bound $U$. Adding any finite number of feasible, correctly evaluated witnesses and replacing $U$ by the minimum of all their upper bounds preserves containment of the original feasible infimum. Every returned witness with $U-L\le\delta$ has original Bellman regret at most $\delta$. The conclusion is independent of the training method and of whether an added witness is a vertex of the lower cover.
\end{proposition}
\begin{proof}
The unchanged cover retains its lower bound. Every evaluated feasible action supplies an upper bound for the infimum; their minimum remains an upper bound. The returned witness has cost at most its own upper endpoint.
\end{proof}

The vector experiment compares adaptive critical-triangle and uniform largest-edge refinement with two-output ReLU and quadratic proposals under the same certificate. Uniform refinement changes the triangle selected for a split, not the global lower-bound calculation. Four-date neural proposals transfer three-date fitted weights. Vector fitting uses action-label loss and is not relabeled as the scalar certificate-aware training theorem. The complete development companion gives the full training, validation, reuse and witness specification.

\section{Two controls and two continuous innovations}\label{sec:vector67}
The absence of a state lattice does not remove action or innovation dimensions. We extend the query implementation to the previously specified two-control economy of Section~\ref{sec:accuracy62}: $a_1,a_2\ge0$, $a_1+a_2\le c(x)$, the interaction cost is $a_1a_2/4$, and the second innovation is an independent common-direction uniform shock. The original scalar model and its experiments are unchanged.

At a queried state, triangulate the feasible action simplex. Let $v_0,v_1,v_2$ be a triangle's vertices with verified lower endpoints $l_0,l_1,l_2$. If $\theta_j(a)$ are its barycentric coordinates, define
\begin{equation}\label{eq:triangle67}
 p(a)=\sum_{j=0}^2\theta_j(a)l_j
       -\frac H2\sum_{j=0}^2\theta_j(a)\|v_j-a\|_2^2.
\end{equation}
This is the quadratic $H\|a\|_2^2/2+b^\top a+c$ interpolating the lower endpoints. Its approximate minimizer is only a proposed support point. For any point $y\in\mathbb R^2$, including a numerically displaced point outside the triangle,
\begin{equation}\label{eq:dual-triangle67}
 p(y)+\min_j\nabla p(y)^\top(v_j-y)
                    \le \min_{a\in\operatorname{co}(v_0,v_1,v_2)}p(a).
\end{equation}
Thus an outward supporting-plane calculation, not an uncertified quadratic solver, supplies the lower endpoint.

\begin{theorem}[Simplicial Bellman query]\label{thm:vector67}
For the two-control primitives, \eqref{eq:triangle67} is a lower minorant of the true Bellman objective. Taking the smallest valid triangle lower endpoint and the best feasible-action upper endpoint gives a lower--upper certificate for the full continuous simplex. Include the omitted capacity strip and actual action-rounding allowances. Returning when this gap is at most $\delta$ gives the same acquired-policy loss recursion as Theorem~\ref{thm:query63}.

A failed floating triangulation can be replaced by an exact-rational simplex enumeration. Under the stated positive precision slack, this recovery is finite at every finite horizon. With two innovations divided into $q$ bins each, a valid local quadrature error is
\begin{equation}\label{eq:two-shock67}
 \zeta=\frac{5\beta M_{t+1}d}{24576q^2}.
\end{equation}
If $A_r$ bounds the number of action queries at a state with $r$ dates remaining, a work account has the form
\begin{equation}\label{eq:vector-work67}
 W_r\le C_d A_r+A_rq_r^2W_{r-1},\qquad C_d=O(d),
\end{equation}
with separate charges for triangulation, arithmetic and output. There is no tensor state-table factor, but the descendant tree remains.
\end{theorem}
\noindent The proof is reproduced in Supplement, Section~\ref{supp:relocated69-4}.

The executed algorithm uses longest-edge refinement of the lowest triangle and an outward terminal convexity-gap check. Its floating cap is followed by the stated recovery rule, not by an uncertified returned action. The implementation batches descendants but can allocate a descendant array before slicing it into query batches; a reported maximum query batch is therefore not a bound on total live state memory. We report process memory separately. The preceding two-control execution is a verified extension of the NBO query backend, not evidence for a trained two-output selector. The new comparison in Section~\ref{sec:study69} separately evaluates trained two-output proposals under this unchanged certificate.

\section{From returned policies to economic resource decisions}\label{sec:decisions49}
An all-state upper bound and a comparison of actual policy costs answer different economic questions. We now connect the policies returned by a construction frontier to direct cost inference and then to a choice of implementation resources. The construction data select the policies; independently generated evaluation paths assess them. This separation permits a valid comparison of first-crossing policies with different resolutions.

\subsection{Simultaneous inference after construction}
Let $\mathcal H$ denote all construction information: primitive inputs, numerical labels, fitted continuations, returned policies and complete certificate ladders. Conditional on $\mathcal H$, fix a finite catalogue of implemented-policy contrasts. Each implementation includes its observation rule, original witness or nearest-node selector, true-state feasible repair and action quantization. For a contrast $i$, let $D_i$ be the difference in expected discounted costs under a specified initial law. A common-innovation path score evaluates the two policies on their own endogenous trajectories. It is not the difference between two regret bounds.

Suppose the initial and innovation uniforms are partitioned into equal bins. Independent uniform bin indices are drawn, and interval simulation encloses the score for every continuous realization inside the drawn bins. All unresolved state, selector and repair branches remain in the enclosure. If $[A_i,B_i]$ is a deterministic support interval, clip the path endpoints to that support. Write $\underline m_i,\overline m_i$ for outward sample-mean bounds on the lower and upper endpoint samples and $v_{i,L},v_{i,U}$ for upper bounds on their unbiased sample variances. For $n\ge2$ define
\begin{equation}\label{eq:radius49}
 R_i(v)=\sqrt{2v\ell/n}+\frac{7(B_i-A_i)\ell}{3(n-1)},
 \qquad \ell\ge\log(4K/\alpha),
\end{equation}
where $K$ bounds the number of sampled contrasts.

\begin{theorem}[Construction-conditional economic intervals]\label{thm:conditional49}
Under the stated independent-bin model, conditional on any valid construction history $\mathcal H$, with probability at least $1-\alpha$ all catalogue contrasts satisfy
\begin{equation}\label{eq:direct49}
 D_i\in[l_i,u_i]:=
 [\max\{A_i,\underline m_i-R_i(v_{i,L})\},
  \min\{B_i,\overline m_i+R_i(v_{i,U})\}].
\end{equation}
The same coverage holds unconditionally. Subsequent selection among these covered contrasts, or evaluation at any deterministic implementation-price vector, needs no additional union bound. This conclusion applies only to policies and initial laws included in the catalogue; it does not create direct-cost observations for every unstudied tolerance or initial state.
\end{theorem}
\noindent The proof is in Supplement, Section~\ref{supp:proof-decisions49-1}.


In particular, selecting the first policy to meet a specified certificate is permitted when that selection uses construction records rather than the subsequent evaluation paths. Repeated construction clocks are not independent draws from a neural-optimizer population. The statistical randomness here belongs to the evaluation design. A fixed pseudorandom stream implements a reproducible experiment; it is not itself a proof of the independent-bin assumption.

\subsection{Replacement-fee identification}
Write $D=J^W-J^F$. If witness is installed, replacing it by FVI changes total cost by $k-D$, where $k\ge0$ is the replacement fee. If FVI is installed, the change is $k+D$. Thus the two directional gains before fees are $D$ and $-D$.

\begin{corollary}[The break-even fee frontier]\label{cor:fee49}
On an event where $D\in[l,u]$, the two directional gains belong to $[l,u]$ and $[-u,-l]$. The unknown absolute cost difference satisfies
\begin{equation}\label{eq:absolute49}
 |D|\in[\operatorname{dist}(0,[l,u]),\max\{|l|,|u|\}].
\end{equation}
For every $k\ge\max\{0,u,-l\}$ neither replacement strictly lowers expected cost. By contrast, $k\ge\max\{0,l,-u\}$ only removes a \emph{certified} strictly profitable replacement. The two thresholds generally differ. Strict inequalities in the first criterion certify that both replacements strictly increase cost.
\end{corollary}
\noindent The proof is in Supplement, Section~\ref{supp:proof-decisions49-2}.


The fee frontier is more informative than a conclusion at one generous charge. The earlier $1/64$ replacement decision is retained with its original interpretation and confidence account. The new analysis reports the two directional gain intervals and both thresholds instead of selecting a favorable fee after seeing the data. There is no welfare calibration hidden in the units: a decision maker must supply the economically relevant charge, or use the full frontier to see what remains unidentified.

\subsection{Observation costs on the same bit grid}
Let $\mathcal B=\{4,\ldots,16\}$ be the available bit precisions. For one fixed returned controller, write $J_b$ for the expected cost of its actual $b$-bit, robustly repaired and quantized implementation. Use $b_0=16$ as a common reference and obtain simultaneous intervals for $D_b=J_b-J_{b_0}$. The identity interval at $b=b_0$ is exactly $[0,0]$. Let $s_b=d b\sum_{t<T}B_t$ be the discounted number of coordinate-bits purchased; the net cost is $J_b+\lambda s_b$.

\begin{theorem}[Certified choice of actual net implementation cost]\label{thm:net49}
Suppose all $D_b\in[l_b,u_b]$ simultaneously. For any $\lambda\ge0$ define
\begin{align}\label{eq:net49}
 L(\lambda)&=\min_b\{l_b+\lambda s_b\},&
 U(\lambda)&=\min_b\{u_b+\lambda s_b\},\\
 \widehat b(\lambda)&=\min\arg\min_b\{u_b+\lambda s_b\}.\notag
\end{align}
Then, simultaneously for all $\lambda\ge0$,
\begin{equation}\label{eq:netregret49}
 0\le J_{\widehat b}+\lambda s_{\widehat b}
       -\min_{b\in\mathcal B}(J_b+\lambda s_b)
 \le U(\lambda)-L(\lambda).
\end{equation}
The lower and upper functions are finite piecewise-affine envelopes. All changes in their active lines occur at nonnegative intersections of catalogue lines, with exact comparisons deciding ties. The smallest selected precision $\widehat b(\lambda)$ is weakly decreasing in $\lambda$.
\end{theorem}
\noindent The proof is in Supplement, Section~\ref{supp:proof-decisions49-3}.


This result makes a statement about \emph{actual} net cost within the declared bit catalogue. It differs from minimizing a uniform all-state loss bound. Both criteria are useful: the latter gives distribution-free protection relative to the full-information optimum, while \eqref{eq:netregret49} uses a specified initial law to compare feasible implementations of one returned policy. Neither criterion gives an optimum over arbitrary sensor technologies or all policies.

For a fair comparison the conventional actor receives its own acquisition allowance. Suppose its nearest-node policy has a certified selected-policy account at a reported state, obtained from the graph modulus $L_t^{\mathrm{gr}}$, and its actual interpolated critic has modulus $L_t^f$. Transporting the same selected node action to the true state, then using cell-wide capacity repair, adds at most
\begin{equation}\label{eq:fvisensor49}
 (L_t^{\mathrm{gr}}+L_t^f+2D_tK_t^A)r(b)+D_t\Delta_t.
\end{equation}
Indeed, transport the feasible node pair directly to the true state, bound its node distance by the reported-state distance plus $r(b)$, and bound $f_t(\widehat x)-f_t(x)$ by $L_t^fr(b)$. Robust repair adds $D_t(2K_t^Ar(b)+\Delta_t)$. This proves \eqref{eq:fvisensor49} without assuming actor continuity. For a witness critic $L_t^f=L_t^{\mathrm{gr}}$, it reduces to the preceding witness allowance. The empirical comparison therefore evaluates the all-state-account minimizer, the interval upper-net-cost selector, and a descriptive midpoint selector on the \emph{same} integer bit and price grid. The midpoint selector is not called the true optimum.

\section{Policy-cost inference at the scale of verified regret}\label{sec:centered69}
The query certificate and a direct cost experiment answer different questions. Nevertheless, a certificate can improve the experiment without replacing its estimand. The relevant observation is a sum of optimal Bellman advantages along the actual policy trajectory. Its expectation is policy loss, while its range is controlled by the local accuracy contract rather than by the range of economic costs. All conditional expectations used to construct this observation must themselves be evaluated and charged.

Write $V_t$ for the original Bellman optimum, $V_T=g$, $b_0=1$, and $b_t=\prod_{j<t}\beta_j$. For a feasible frozen policy $\pi$, define
\[
 a_t^\pi(x)=c_t(x,\pi_t(x))+\beta_tP_t^{\pi_t(x)}V_{t+1}(x)-V_t(x),
 \qquad Z^\pi=\sum_{t<T}b_t a_t^\pi(X_t^\pi).
\]
The symbol $a_t^\pi$ denotes an advantage, not an action. Assume a verified local bound $0\le a_t^\pi(x)\le e_t^\pi$ for every true state, including the acquisition allowance. Put $G_\pi=\sum_{t<T}b_t e_t^\pi$. A sample contains the actual trajectory of each compared controller under the original law; common initial states and common innovations may couple the controllers. Different sample rows are independent conditional on the frozen construction.

\begin{theorem}[Bellman-centered cost comparison]\label{thm:centered69}
For two such controllers $\pi,\nu$ with a common initial law $\lambda$, let $D=Z^\pi-Z^\nu$. Then
\begin{equation}\label{eq:centered-identity69}
 \mathbb E D=\mathbb E_\lambda[J_0^\pi-J_0^\nu],
 \qquad -G_\nu\le D\le G_\pi.
\end{equation}
Suppose deterministic computation encloses $D_i$ by $[L_i,U_i]\subseteq[-G_\nu,G_\pi]$ on every sampled continuous bin, including every compatible acquired action. For a conditionally fixed finite family of $M$ contrasts, put $R=G_\pi+G_\nu$ and
\[
 h_n=R\sqrt{\frac{\log(2M/\alpha)}{2n}}.
\]
With probability at least $1-\alpha$, simultaneously for the family,
\begin{equation}\label{eq:centered-ci69}
 \mathbb E_\lambda[J_0^\pi-J_0^\nu]
 \in\left[\frac1n\sum_iL_i-h_n,\frac1n\sum_iU_i+h_n\right].
\end{equation}
The interval may be intersected with $[-G_\nu,G_\pi]$. Its width is the mean numerical enclosure width plus $2h_n$. More generally, any valid endpoint-sample concentration bound may replace the bounded-sum radius.
\end{theorem}
\begin{proof}
Conditional expectation and $b_{t+1}=b_t\beta_t$ imply
\[
 \mathbb E\sum_{t<T}b_t a_t^\pi(X_t^\pi)
 =\mathbb E\left[\sum_{t<T}b_tc_t(X_t^\pi,\pi_t(X_t^\pi))
                  +b_Tg(X_T^\pi)-V_0(X_0^\pi)\right].
\]
The intermediate value terms cancel. Subtracting the two expectations cancels the initial value under the common initial law, proving the identity. Summing the pointwise local bounds proves the range. Applying the bounded-sum inequality to each lower and upper endpoint sample, with failure allocation $\alpha/(2M)$, and taking a union bound proves \eqref{eq:centered-ci69}. Endpoint containment supplies the comparison between endpoint expectations and the target. The width identity follows by subtraction. No pathwise identity between $D$ and the realized cumulative-cost difference is asserted.
\end{proof}

The construction is related to martingale control variates and sequential value-based evaluation, including \citet{jiang2016}. It is not an off-policy importance-weighting argument and does not assume that a fitted value is correct. The additional requirements here are an original-optimum advantage enclosure, a true-state local regret bound, and a paid numerical evaluation of the conditional value. The target in \eqref{eq:centered-ci69} is the original expected cost difference, not the difference between two regret certificates. The realized centered observation generally differs from the realized cash-cost difference; their expectations agree by the displayed identity.

\subsection{Numerical construction and a target-precision rule}
The returned action is fixed before evaluation. At a center $\bar x$ of a true-state box, let $[Q_L,Q_U]$ enclose $Q_t^*(\bar x,\pi_t(x))$ for that fixed action and let $[V_L,V_U]$ enclose $V_t(\bar x)$. If $L_Q,L_V$ are state Lipschitz bounds in the $\ell_1$ norm and $\rho$ bounds $\|x-\bar x\|_1$, then
\begin{equation}\label{eq:advantage-box69}
 [Q_L-V_U-(L_Q+L_V)\rho,
   Q_U-V_L+(L_Q+L_V)\rho]\cap[0,e_t^\pi]
\end{equation}
is a valid advantage enclosure. An empty intersection stops verification. Actions must be feasible at every state used by the endpoint queries. If the box crosses an acquisition boundary, one may enumerate all compatible acquired actions, or use $[0,e_t^\pi]$ for the unresolved part. The latter is conservative and its width remains in the reported numerical account. It is not permissible to discard that trajectory.

For the scalar investment economy, at $r$ remaining dates we use
\[
 L_Q=\frac8d+\frac{11\beta}{16}G_{r-1},\qquad L_V=G_r,
 \qquad e_t^\pi=\delta_t+d(L_Q+L_V)2^{-b},
\]
from the original acquisition argument. The evaluator queries the optimum with its own, smaller tolerance; it does not retrain the policy or change the action being evaluated. Original-law quadrature, action-cover error and arithmetic are already included in $[Q_L,Q_U]$ and $[V_L,V_U]$. They must not be omitted or counted again as independent random noise.

\begin{corollary}[Precision and noninferiority]\label{cor:precision69}
Fix a numerical half-width allocation $v>0$ and a sampling allocation $s>0$ before observing the new contrasts. If $n^{-1}\sum_i(U_i-L_i)\le2v$ and
\[
 n\ge\frac{R^2\log(2M/\alpha)}{2s^2},
\]
then the interval in \eqref{eq:centered-ci69} has half-width at most $v+s$. For a prespecified loss margin $m>0$, an upper endpoint at most $m$ establishes noninferiority at that margin; containment in $[-m,m]$ establishes the corresponding two-sided equivalence statement on the same event. Neither statement identifies a strict advantage.
\end{corollary}
\begin{proof}
The sample-size condition implies $h_n\le s$. Insert this and the numerical-width condition into the width identity. The final claims are set inclusions for the simultaneous confidence interval.
\end{proof}

Both $n$ and evaluator precision have costs. The deposited design fixes them prospectively and reports failure to meet the width allocation rather than increasing the sample after inspecting a favorable sign. The actual numerical width is decomposed into the conditional Bellman enclosure, the state-bin Lipschitz expansion, unresolved acquisition branches and endpoint arithmetic. The conditional Bellman component includes its action-cover and continuous-shock quadrature work. A smaller statistical support is therefore not advertised as a free runtime improvement.

\subsection{The adoption decision after inference}
Let $[l,u]$ enclose $\mathbb E[J^\pi-J^\nu]$, let $W_\pi,W_\nu$ be measured complete deployment work under a declared boundary, and let $p\ge0$ be a decision maker's price per unit of that work. Replacing $\nu$ by $\pi$ at fee $f\ge0$ has net-cost difference in
\[
 [l+p(W_\pi-W_\nu)+f,\ u+p(W_\pi-W_\nu)+f].
\]
A negative upper endpoint licenses the replacement under those declared prices. A noninferiority finding permits a separate assessment of computational savings; it does not convert a normalized loss margin into an estimated monetary value. This keeps the original economic purchase question while making its statistical precision and its required price inputs explicit.

\section{Training, policy costs, and deployment in the investment economy}\label{sec:study69}
The economic primitives, feasible set and continuous innovation laws are those of Sections~\ref{sec:accuracy62} and \ref{sec:query63}. The new experiment changes the fitted proposal, verification workload and evaluation precision, not the economy. Every callable controller acquires the state downward and returns only certificate-authorized feasible actions. The hypothesis that matters for deployment is a joint statement about achieved policy cost and the complete work of returning that policy.

The preceding comparison is retained in full in the complete development companion, with its original records and inference families. It found adverse cold-start ReLU timings, no ReLU reuse win, only a late and task-specific quadratic crossing, and unresolved cost differences. The original large four-date vector calculation and the corrected cumulative-array counterexample remain visible. New experiments do not overwrite those findings or reinterpret repeated observations as new samples.

\subsection{A prospectively fixed comparison}
The scalar tasks have $(d,T)=(2,2)$ and $(8,3)$, with original-optimum loss targets $1/128$ and $1/64$. They use twenty-four-bit acquired states and a thirty-two-bit downward action lattice. The certificate-aware ReLU and Bernstein readouts, the label-loss ReLU and the quadratic proposal receive identical purchased training information. Two fixed fitting seeds and four separate process repetitions are declared in advance. Adaptive and upper-chord minimization remain prediction-free controls. The earlier linear, spline, nearest-action, derivative, bisection and unrestricted ReLU comparisons retain their original place in the historical catalogue; they are not described as newly rerun arms.

Independent validation purchases 256 fresh contexts at every nonterminal date. The sixty-four-block reuse design fixes new workload streams before deployment, pays for every model load and identity check, and includes the declared stale-identifier fallback. The two-control comparison uses the same coupled-capacity economy and two independent innovations, with adaptive and uniform triangular refinement and two trained two-output proposals. Four-date neural proposals use the explicitly declared three-date transfer.

All planned parent-process receipts, finite optimizer statuses, warning messages, timeouts and return failures belong to the catalogue. Counts of services are not probabilities of successful learning. Model identities, data, exact optimization certificates and output traces bind each policy to its own construction.

\input{tables/findings69}

\subsection{Original policy costs at a declared precision}
The inference experiment compares adaptive search, the two newly certificate-trained readouts, and the original frozen ReLU callable controller from the preceding study. The last arm loads its original model rather than retraining it. There are 1,024 new common-path rows per task under the original continuous initial and innovation laws. Forty-eight-bit continuous bins enclose the sampled paths. The finer evaluator tolerance is fixed at one thirty-second of the policy's declared target.

The centered observations in Theorem~\ref{thm:centered69} have the original policy-cost difference as their expectation. Their smaller support follows from verified local regret, while their numerical enclosure requires additional original Bellman evaluations. The ordinary cumulative cash-cost estimator is computed on exactly the same new paths. The comparison therefore exposes the gain in resolving power and the additional evaluation work, rather than confusing a regret bound with a measured economic difference.

\input{tables/centered69}

The two-sided design margin is $1/256$ in normalized economic-cost units. It was fixed before the new contrasts, with numerical and sampling half-width allocations of $1/512$ each. The supplementary precision table reports whether each allocation was achieved. Equivalence is asserted only for intervals contained in the declared margin, and noninferiority only when the relevant upper endpoint satisfies it. These findings do not identify a strict policy gain and do not calibrate the normalized margin as money. A negative or positive empirical mean alone licenses neither adoption nor rejection.

The numerical ledger distinguishes the conditional Bellman interval, the state-box expansion, unresolved acquisition branches and final endpoint arithmetic. The state-box expansion combines continuous-bin uncertainty with propagated state arithmetic; those two components are not separately identified in the implementation. Intersections with valid local-regret bounds can reduce the final width, so raw component widths before intersection need not add to the printed final interval. Ambiguous acquisition branches are conservatively retained, not omitted from the sample. Combining the separately allocated validation, centered-cost and cash-cost families requires their union bound.

\subsection{Achieved training and complete cold-start work}
The supplementary readout table reports the achieved convex objective and exact box dual gap at the returned quantized weights. The validation table compares the work-loss account with actual full-action recovery queries on predictor-independent contexts. A validated root-query upper bound is not a bound on process seconds, and a valid policy return is not evidence that a hidden-layer optimizer converges reliably.

\input{tables/cold69}

The complete cold-start clock begins before process launch and includes imports, purchased training information, fitting and exact readout certification, prediction, all Bellman descendants, path enclosure, serialization, durable output and shutdown. Parent receipt bookkeeping is additionally charged in the cohort total. Machine-independent query counts, rational recovery, training queries and largest live batches remain separate fields. The timing range includes every declared repetition. Available frequency and governor information is recorded, but frequency is not claimed experimentally locked; close clock differences are consequently descriptive rather than identified hardware-invariant gains.

\subsection{Reuse after paying for the model}
\input{tables/reuse69}
A final complete-process win and an early persistent prefix crossing are different events. The former includes shutdown and the final output account; the latter requires a crossing by block forty-eight that survives every later block. A model can fail either criterion, and both outcomes remain in the table. The sixty-four-block design extends the original prospective window without treating the absence of a later reversal as evidence about unobserved future blocks. Reuse adoption must combine these measured costs with the policy-cost interval and the decision maker's explicit computation price and replacement fee.

\subsection{Two controls and horizon dependence}
\input{tables/vector69}
The additional vector upper witness and the unchanged triangular lower cover implement Proposition~\ref{prop:vector-witness69}. Both learned arms use the same paid label states. The comparison distinguishes the effects of a proposal from the effects of uniform or adaptive action refinement. It is not a derivative branch-and-bound comparison, and vector label fitting is not relabeled as certificate-aware scalar training.

The horizon-four record includes memory, floating and rational query counts, training or transfer identity and complete process time. The previously observed large descendant-tree cost remains a warning about horizon and innovation dependence even when the state lattice is absent. These finite path diagnostics are not a simultaneous population policy-cost comparison for vector controllers. The independent rational implementation checks every represented new mid-bin trajectory and independent scalar and two-control terminal integrals; it does not independently re-integrate every node of the largest inherited simplicial tree.

The resulting economic distinction is precise. The method returns an implementable policy under an original-optimum contract; an achieved readout certificate measures what the fitting procedure actually produced; a prospective work account prices its use; and centered evaluation can assess policy differences at a prespecified scale. None of those objects is substituted for the others. The original NBO construction and its broader applications retain their stated hypotheses, while the new numerical conclusions concern the particular expected-cost economies evaluated here.


\section{Conclusion}\label{sec:conclusion67}
Neural Bellman Operators connect continuation construction to feasible economic decisions under explicit resource and accuracy accounts. The construction does not need a successful nonconvex fitting theorem: arbitrary proposals remain subject to an original-law certificate and a finite recovery procedure. The fixed-feature readout adds a separately computable optimization certificate, and its achieved work loss can be assessed on an independent context law.

Bellman-centered inference supplies a further connection. The original policy-cost difference is the expectation of a regret-scale observation whose numerical evaluation is itself verified and charged. This permits a precision-based comparison rather than an attempted policy ranking from two approximation bounds. Noninferiority at a declared normalized margin is different from strict advantage, a measured computation saving, or an estimated monetary welfare gain. A purchase decision must combine the relevant interval with explicit computation prices and replacement costs.

The complete numerical catalogue retains unfavorable as well as favorable findings. Removing a state lattice does not remove horizon or innovation-tree costs, and the optional vector witness does not alter that fact. The resulting theorem--implementation--cost chain remains part of the original NBO paper. Its broader applications are preserved without being presented as numerical validations of the specific expected-cost theorem established here.

\bibliographystyle{ecta-fullname}
\bibliography{references}
\end{document}
